\cotefile{157-5}
\scanpage{85}
\mspage{72}

alors l'axiome de factorisation \inserted{$\mathrm{Cof}_3$} est satisfait. Inversement, si $\mathrm{Cof}_3$ est satisfait, et si $W$ est stable par facteurs directs, alors on a 4.4. Donc si $W$ est stable par facteurs directs, $\mathrm{Cof}_3$ équivaut à (4.4.)

4.2. Il me semble que ces conditions, en pratique, ne seront guère jamais vérifiées, \unsure{que|car} si $W$ (qui est donné en \deleted{tout} premier lieu,\inserted{\unsure{satisfaisant à} $\underline{\mathrm{Cof}_0^{*}}$}) \margin{$*$ i.e. \unsure{saturation|stabilité} faible}
% slanted
n'est \deleted{\unsure{satisfai}} pas stable par $\unsure{\varinjlim|\underline{\lim}}$ filtrantes.

À ce moment, \unsure{pour} construire la structure à cofibrations, on part de $C_0$ \unsure{petite|petites} partie de $\operatorname{Fl}M$, \unsure{disons|choisie} telle que l'on ait
\[
(C_0)_{*}\subset W,\tag*{(4.5.)}
\]
on \unsure{prend|peut} pour \inserted{au choix} $C$ le saturé \inserted{$C'_0$} de $C_0$ par cochg\textsuperscript{t}. de base et composition transfinie,, \deleted{\unsure{ou}} $C=\widetilde C_0={}$ ens des facteurs directs de $C'_0$. Les conditions $\mathrm{Cof}_1$,\inserted{(stabilité de $C$)} $\mathrm{Cof}_3$\inserted{(factorisation)} sont satisfaites automatiquement, et il reste à s'occuper de $\mathrm{Cof}_2$, qui revient à la stabilité de $C\cap W$ par cochg\textsuperscript{t}. de base (la stabilité par comp., et que ça contienne les iso, étant \unsure{immédiat|immédiate}). C'est une condition délicate, et la seule façon \inserted{\unsure{standard|standart}} que je vois pour l'assurer, c'est en supposant
\[
C_0\subset\mathrm{Cof}_W\,;\tag*{(4.6.)}
\]
% archive sheet number: 84
